\begin{itemize}
\item Vandermonde's Identity
$$
C(n + m, k) = \sum_{i=0}^k C(n, i)C(m, k - i)
$$

\item Kirchhoff's Theorem

Denote $L$ be a $n \times n$ matrix as the Laplacian matrix of graph $G$, where $L_{ii} = d(i)$ (in-degree for directed case), $L_{ij} = -c$ where $c$ is the number of edge $(i, j)$ in $G$.
\begin{itemize}
    %\itemsep-0.5em
    \item The number of undirected spanning trees in $G$ is $\lvert \det(\tilde{L}_{11}) \rvert$.
    \item The number of directed spanning tree rooted at $r$ in $G$ is $\lvert \det(\tilde{L}_{rr}) \rvert$.
\end{itemize}

\item BEST Theorem

The number of Eulerian circuits of a directed Eulerian graph is $t_w(G) \prod_v (\deg^+(v) - 1)!$, where $t_w(G)$ is the number of spanning trees rooted at $w$ (Kirchhoff; the same for every $w$).

\item Tutte's Matrix

Let $D$ be a $n \times n$ matrix with $d_{ij} = x_{ij}$ ($x_{ij}$ chosen uniformly at random) if $i < j$ and $(i, j) \in E$, $d_{ij} = -d_{ji}$ if $i > j$, and $0$ otherwise. $\frac{rank(D)}{2}$ is the maximum matching on $G$ (compute the rank modulo a large prime; correct with high probability).

\item Tutte–Berge Formula

Maximum matching $= \frac12 \min_{U \subseteq V} \big(|V| + |U| - \mathrm{odd}(G - U)\big)$, where $\mathrm{odd}$ counts the components with an odd number of vertices.

\item Lindström–Gessel–Viennot Lemma

In a DAG with sources $a_1, \ldots, a_n$ and sinks $b_1, \ldots, b_n$, let $e(a, b)$ be the number of paths $a \to b$ (or the sum of their weights) and $M_{ij} = e(a_i, b_j)$. Then $\det M = \sum_{\sigma} \operatorname{sgn}(\sigma) \cdot \#\{\text{vertex-disjoint paths } a_i \to b_{\sigma(i)}\}$. If only $\sigma = \mathrm{id}$ can be disjoint (e.g. a grid with sources and sinks in the same order), $\det M$ counts the non-intersecting path systems.

\item Cayley's Formula

\begin{itemize}
    %\itemsep-0.5em
  \item Given a degree sequence $d_1, d_2, \ldots, d_n$ for each \textit{labeled} vertices, there are $\frac{(n - 2)!}{(d_1 - 1)!(d_2 - 1)!\cdots(d_n - 1)!}$ spanning trees.
  \item Let $T_{n, k}$ be the number of \textit{labeled} forests on $n$ vertices with $k$ components, such that vertex $1, 2, \ldots, k$ belong to different components. Then $T_{n, k} = kn^{n - k - 1}$.
\end{itemize}

\item Cycle Lemma

If integers $x_1, \ldots, x_N$ are all $\le 1$ and sum to $k > 0$, exactly $k$ of the $N$ cyclic shifts have every prefix sum $> 0$. Grouping all permutations of these numbers by rotation, exactly a $\frac{k}{N}$ fraction of them are good.
\begin{itemize}
  \item Ballot: A gets $a$ votes, B gets $b$ votes. Orders in which, at every moment, A has strictly more than $t$ times B's votes: $\frac{a - tb}{a + b}\binom{a + b}{a}$ (write A as $+1$, B as $-t$); at least $t$ times: $\frac{a + 1 - tb}{a + 1}\binom{a + b}{a}$ (count strict orders with $a + 1$ A votes, then drop the first vote, which is always A).
  \item Rooted unlabeled trees (whose children have a left-to-right order): two trees are equal iff their roots have the same number of children and all of their $i$-th subtrees are equal (defined recursively). Given $N = \sum n_d$ nodes, $n_d$ of which have $d$ children, write $1 - (\#\text{children})$ for each node; the sequences whose prefix sums are all $\ge 1$ are exactly the post-order traversals. The number of distinct trees is $\frac{1}{N}\binom{N}{n_0, n_1, \ldots}$.
\end{itemize}

\item Erdős–Gallai theorem

A sequence of nonnegative integers $d_1\ge\cdots\ge d_n$ can be represented as the degree sequence of a finite simple graph on $n$ vertices if and only if $d_1+\cdots+d_n$ is even and $\displaystyle\sum_{i=1}^kd_i\le k(k-1)+\displaystyle\sum_{i=k+1}^n\min(d_i,k)$ holds for every $1\le k\le n$.

\item Gale–Ryser theorem

A pair of sequences of nonnegative integers $a_1\ge\cdots\ge a_n$ and $b_1,\ldots,b_n$ is bigraphic if and only if $\displaystyle\sum_{i=1}^n a_i=\displaystyle\sum_{i=1}^n b_i$ and $\displaystyle\sum_{i=1}^k a_i\le \displaystyle\sum_{i=1}^n\min(b_i,k)$ holds for every $1\le k\le n$.

\item Fulkerson–Chen–Anstee theorem

A sequence $(a_1,b_1),\ldots,(a_n,b_n)$ of nonnegative integer pairs with $a_1\ge\cdots\ge a_n$ is digraphic if and only if $\displaystyle\sum_{i=1}^n a_i=\displaystyle\sum_{i=1}^n b_i$ and $\displaystyle\sum_{i=1}^k a_i\le \displaystyle\sum_{i=1}^k\min(b_i,k-1)+\displaystyle\sum_{i=k+1}^n\min(b_i,k)$ holds for every $1\le k\le n$.

\item Dilworth / Mirsky

\begin{itemize}
  \item Minimum number of chains covering a poset $=$ maximum antichain; minimum number of antichains $=$ longest chain.
  \item Minimum vertex-disjoint path cover of a DAG $= n\, -$ maximum matching (split each vertex into out/in). If paths may share vertices, take the transitive closure first.
\end{itemize}

\item Graph facts

\begin{itemize}
  \item Hall: a bipartite graph has a matching saturating $L$ iff $|N(S)| \ge |S|$ for all $S \subseteq L$; maximum matching $= |L| - \max_{S \subseteq L} (|S| - |N(S)|)$.
  \item Gallai: maximum independent set $+$ minimum vertex cover $= n$; with no isolated vertices, maximum matching $+$ minimum edge cover $= n$.
  \item Dirac / Ore ($n \ge 3$): a Hamiltonian cycle exists if every $\deg(v) \ge n/2$, or if $\deg(u) + \deg(v) \ge n$ for all non-adjacent $u, v$.
\end{itemize}

\item Burnside's Lemma

Number of orbits $= \frac{1}{|G|} \sum_{g \in G} |X^g|$, $X^g$: elements fixed by $g$.
\begin{itemize}
  \item Necklaces ($n$ beads, $k$ colours, rotations): $N = \frac1n \sum_{d \mid n} \varphi(d)\, k^{n/d}$.
  \item Bracelets (also reflections): $\frac N2 + \frac12 k^{(n+1)/2}$ for odd $n$; $\frac N2 + \frac14 (k+1) k^{n/2}$ for even $n$.
\end{itemize}

\item Labeled graph counts

\begin{itemize}
  \item Connected graphs: EGF $C(x) = \ln G(x)$, $G(x) = \sum_n 2^{\binom n2} \frac{x^n}{n!}$ (use \texttt{Ln}).
  \item DAGs: $a_n = \sum_{k=1}^{n} (-1)^{k+1} \binom nk 2^{k(n-k)} a_{n-k}$, $a_0 = 1$.
\end{itemize}

\item Hook length formula

Number of standard Young tableaux of shape $\lambda$ ($n$ cells) is $\frac{n!}{\prod_{c \in \lambda} h(c)}$, where $h(c)$ = cells to the right of $c$ + cells below $c$ + 1.

\item Fibonacci

\begin{itemize}
  \item $M = \begin{bmatrix}1&1\\1&0\end{bmatrix}$, $M^n = \begin{bmatrix}F_{n+1}&F_n\\F_n&F_{n-1}\end{bmatrix}$; the next line follows from $M^{m+n} = M^m M^n$ and $\det M^n = (-1)^n$.
  \item $F_{m+n} = F_m F_{n+1} + F_{m-1} F_n$, \quad $F_{n-1}F_{n+1} - F_n^2 = (-1)^n$
  \item $\gcd(F_m, F_n) = F_{\gcd(m,n)}$; $F_m \mid F_n \iff m \mid n$ ($m \ge 3$). Same shape: $\gcd(a^m - 1, a^n - 1) = a^{\gcd(m,n)} - 1$.
  \item Pisano period $\pi(m) \le 6m$. Zeckendorf: every $n$ is uniquely a sum of non-adjacent Fibonacci numbers (greedy).
\end{itemize}

\item Generalized Euler

For any $a$ and $b \ge \varphi(m)$: $a^b \equiv a^{\,b \bmod \varphi(m) + \varphi(m)} \pmod m$ (the condition on $b$ is required).

\item Landau's theorem

$s_1 \le \dots \le s_n$ is the score sequence of a tournament iff $\sum_{i \le k} s_i \ge \binom k2$ for all $k$, with equality at $k = n$.

\item Mis\`ere Nim

The first player wins iff (some pile $> 1$ and $\bigoplus a_i \ne 0$) or (all piles $\le 1$ and $\bigoplus a_i = 0$).

\item Pick's theorem

For simple polygon, when points are all integer, we have $A=\text{\#\{lattice points in the interior\}} + \frac{\text{\#\{lattice points on the boundary\}}}{2} - 1$.

\item Möbius inversion formula

\begin{itemize}
    %\itemsep-0.5em
  \item $f(n)=\sum_{d\mid n}g(d)\Leftrightarrow g(n)=\sum_{d\mid n}\mu(d)f(\frac{n}{d})$
  \item $f(n)=\sum_{n\mid d}g(d)\Leftrightarrow g(n)=\sum_{n\mid d}\mu(\frac{d}{n})f(d)$
\end{itemize}

\item Spherical cap

\begin{itemize}
    %\itemsep-0.5em
  \item A portion of a sphere cut off by a plane.
  \item $r$: sphere radius, $a$: radius of the base of the cap, $h$: height of the cap, $\cos\theta = (r-h)/r$.
  \item Volume $=\pi h^2(3r-h)/3=\pi h(3a^2+h^2)/6=\pi r^3(2+\cos\theta)(1-\cos\theta)^2/3$.
  \item Area $=2\pi rh=\pi(a^2+h^2)=2\pi r^2(1-\cos\theta)$.
\end{itemize}

\item Lagrange multiplier

Minimise $\sum f_i(x_i)$ subject to $\sum x_i = C$ with every $f_i$ convex: at the optimum $f_i'(x_i) = \lambda$ for all interior $x_i$ (clamp at the bounds); binary search $\lambda$.

\item Nearest points of two skew lines

\begin{itemize}
\item $\text{Line 1}: \boldsymbol{v}_1 = \boldsymbol{p}_1 + t_1\boldsymbol{d}_1$
\item $\text{Line 2}: \boldsymbol{v}_2 = \boldsymbol{p}_2 + t_2\boldsymbol{d}_2$
\item $\boldsymbol{n} = \boldsymbol{d}_1\times \boldsymbol{d}_2$
\item $\boldsymbol{n}_1 = \boldsymbol{d}_1 \times \boldsymbol{n}$
\item $\boldsymbol{n}_2 = \boldsymbol{d}_2 \times \boldsymbol{n}$
\item $\boldsymbol{c}_1 = \boldsymbol{p}_1 + \frac{(\boldsymbol{p}_2 - \boldsymbol{p}_1)\cdot\boldsymbol{n}_2}{\boldsymbol{d}_1\cdot\boldsymbol{n}_2}\boldsymbol{d}_1$
\item $\boldsymbol{c}_2 = \boldsymbol{p}_2 + \frac{(\boldsymbol{p}_1 - \boldsymbol{p}_2)\cdot\boldsymbol{n}_1}{\boldsymbol{d}_2\cdot\boldsymbol{n}_1}\boldsymbol{d}_2$
\end{itemize}

\item Derivatives/Integrals

Integration by parts:
\(\int_a^bf(x)g(x)dx = [F(x)g(x)]_a^b-\int_a^bF(x)g'(x)dx\)
{
  \setlength{\tabcolsep}{1pt}
  \setlength{\columnsep}{0pt}

  \noindent
  \begin{tabular}{|*{20}{>{$\displaystyle}c<{$}|}}
    \frac{d}{dx}\sin^{-1} x = \frac{1}{\sqrt{1-x^2}}
    &
    \frac{d}{dx}\cos^{-1} x = -\frac{1}{\sqrt{1-x^2}}
    &
    \frac{d}{dx}\tan^{-1} x = \frac{1}{1+x^2}
    \\
    \frac{d}{dx}\tan x = 1+\tan^2 x
    &
    \int\tan ax = -\frac{\ln|\cos ax|}{a}
    &
    % \int x\sin ax = \frac{\sin ax-ax \cos ax}{a^2}
    % &
    \\
    \int e^{-x^2} = \frac{\sqrt \pi}{2} \text{erf}(x)
    &
    \int xe^{ax} dx = \frac{e^{ax}}{a^2}(ax-1)
    \\
  \end{tabular}

  \(
    \displaystyle
    \int \sqrt{a^2 + x^2} = \frac{1}{2} \left(x\sqrt{a^2+x^2} + a^2 \operatorname{asinh}(x/a) \right)
  \)
}

\item Min-Max Inclusion Exclusion Principle

(Works in expect values version, i.e., $E[\min]$, too)
$$
\operatorname*{kthmax}_{i \in S} x_i = \sum_{T \subset S} (-1)^{|T|-k} \binom{|T|-1}{k-1} \min_{j \in T} x_j
$$
$$
\operatorname*{kthmin}_{i \in S} x_i = \sum_{T \subset S} (-1)^{|T|-k} \binom{|T|-1}{k-1} \max_{j \in T} x_j
$$

\item Du's Sieve

\begin{itemize}
  \item compute $S(N)=\sum_{i=1}^N f(i)$, find a function $g$ such that $f*g$ is easy to compute
  \item $\sum_{i=1}^N (f*g)(i)=\sum_{i=1}^N g(i)S(\lfloor\frac{N}{i}\rfloor)$
  \item $\implies g(1)S(N)=\sum_{i=1}^N (f*g)(i)-\sum_{i=2}^N g(i)S(\lfloor\frac{N}{i}\rfloor)$
  \item use sqrt decomposition and memoization (sometimes need to precompute small prefix)
\end{itemize}

\end{itemize}